\documentclass[10pt,a4paper]{amsart}
\usepackage{amsmath,amssymb,amsthm,graphicx}
\usepackage{verbatim}
\newenvironment{reference}{\comment}{\endcomment}
%\newenvironment{reference}{}{}
\newenvironment{slogan}{\comment}{\endcomment}
\newenvironment{history}{\comment}{\endcomment}
% For commutative diagrams we use Xy-pic
\usepackage[all]{xy}

% We use 2cell for 2-commutative diagrams.
\xyoption{2cell}
% Theorem environments.
%
\theoremstyle{plain}
\newtheorem{theorem}[subsection]{Theorem}
\newtheorem{proposition}[subsection]{Proposition}
\newtheorem{lemma}[subsection]{Lemma}

\theoremstyle{definition}
\newtheorem{definition}[subsection]{Definition}
\newtheorem{example}[subsection]{Example}
\newtheorem{exercise}[subsection]{Exercise}
\newtheorem{situation}[subsection]{Situation}

\theoremstyle{remark}
\newtheorem{remark}[subsection]{Remark}
\newtheorem{remarks}[subsection]{Remarks}

\numberwithin{equation}{subsection}

% Macros
%
\def\lim{\mathop{\mathrm{lim}}\nolimits}
\def\colim{\mathop{\mathrm{colim}}\nolimits}
\def\Spec{\mathop{\mathrm{Spec}}}
\def\Hom{\mathop{\mathrm{Hom}}\nolimits}
\def\Ext{\mathop{\mathrm{Ext}}\nolimits}
\def\SheafHom{\mathop{\mathcal{H}\!\mathit{om}}\nolimits}
\def\SheafExt{\mathop{\mathcal{E}\!\mathit{xt}}\nolimits}
\def\Sch{\mathit{Sch}}
\def\Mor{\mathop{\mathrm{Mor}}\nolimits}
\def\Ob{\mathop{\mathrm{Ob}}\nolimits}
\def\Sh{\mathop{\mathit{Sh}}\nolimits}
\def\NL{\mathop{N\!L}\nolimits}
\def\CH{\mathop{\mathrm{CH}}\nolimits}
\def\proetale{{pro\text{-}\acute{e}tale}}
\def\etale{{\acute{e}tale}}
\def\QCoh{\mathit{QCoh}}
\def\Ker{\mathop{\mathrm{Ker}}}
\def\Im{\mathop{\mathrm{Im}}}
\def\Coker{\mathop{\mathrm{Coker}}}
\def\Coim{\mathop{\mathrm{Coim}}}

% Boxtimes
%
\DeclareMathSymbol{\boxtimes}{\mathbin}{AMSa}{"02}

%
% Macros for moduli stacks/spaces
%
\def\QCohstack{\mathcal{QC}\!\mathit{oh}}
\def\Cohstack{\mathcal{C}\!\mathit{oh}}
\def\Spacesstack{\mathcal{S}\!\mathit{paces}}
\def\Quotfunctor{\mathrm{Quot}}
\def\Hilbfunctor{\mathrm{Hilb}}
\def\Curvesstack{\mathcal{C}\!\mathit{urves}}
\def\Polarizedstack{\mathcal{P}\!\mathit{olarized}}
\def\Complexesstack{\mathcal{C}\!\mathit{omplexes}}
% \Pic is the operator that assigns to X its picard group, usage \Pic(X)
% \Picardstack_{X/B} denotes the Picard stack of X over B
% \Picardfunctor_{X/B} denotes the Picard functor of X over B
\def\Pic{\mathop{\mathrm{Pic}}\nolimits}
\def\Picardstack{\mathcal{P}\!\mathit{ic}}
\def\Picardfunctor{\mathrm{Pic}}
\def\Deformationcategory{\mathcal{D}\!\mathit{ef}}
\usepackage{cleveref}
\crefname{originaldefinition}{original definition}{original definitions}
\Crefname{originaldefinition}{Original definition}{Original definitions}
\makeatletter
\newlabel{definition-gysin-homomorphism}{{02T8}{source}{}{sourcepointer-1}{}}
\newlabel{definition-gysin-homomorphism@cref}{{[originaldefinition][0][]02T8}{[1][0][]source}{}{}{}}
\makeatother
\crefname{originallemma}{original lemma}{original lemmas}
\Crefname{originallemma}{Original lemma}{Original lemmas}
\makeatletter
\newlabel{lemma-determine-intersections}{{02U6}{source}{}{sourcepointer-2}{}}
\newlabel{lemma-determine-intersections@cref}{{[originallemma][0][]02U6}{[1][0][]source}{}{}{}}
\makeatother
\crefname{originallemma}{original lemma}{original lemmas}
\Crefname{originallemma}{Original lemma}{Original lemmas}
\makeatletter
\newlabel{lemma-gysin-commutes-cap-c1}{{0B72}{source}{}{sourcepointer-3}{}}
\newlabel{lemma-gysin-commutes-cap-c1@cref}{{[originallemma][0][]0B72}{[1][0][]source}{}{}{}}
\makeatother
\crefname{originallemma}{original lemma}{original lemmas}
\Crefname{originallemma}{Original lemma}{Original lemmas}
\makeatletter
\newlabel{lemma-gysin-factors}{{02TO}{source}{}{sourcepointer-4}{}}
\newlabel{lemma-gysin-factors@cref}{{[originallemma][0][]02TO}{[1][0][]source}{}{}{}}
\makeatother
\crefname{originallemma}{original lemma}{original lemmas}
\Crefname{originallemma}{Original lemma}{Original lemmas}
\makeatletter
\newlabel{lemma-gysin-flat-pullback}{{0B71}{source}{}{sourcepointer-5}{}}
\newlabel{lemma-gysin-flat-pullback@cref}{{[originallemma][0][]0B71}{[1][0][]source}{}{}{}}
\makeatother
\crefname{originalreference}{original reference}{original references}
\Crefname{originalreference}{Original reference}{Original references}
\makeatletter
\newlabel{situation-setup}{{02QL}{source}{}{sourcepointer-6}{}}
\newlabel{situation-setup@cref}{{[originalreference][0][]02QL}{[1][0][]source}{}{}{}}
\makeatother
\title{The cap with a Chern class commutes with the Gysin map}
\author{The Stacks Project Authors}
\date{stacks-project:chow.tex:local-snapshot-20261012; GNU FDL 1.2 or later}
\begin{document}
\maketitle
\noindent\textit{Editorial note.} All delivered mathematical text is the author's own, copied from the pinned Stacks Project chapter file chow.tex; no statement or proof step was rewritten, reordered or removed.\par\medskip
\noindent\textit{Editorial note.} Preamble provenance: the pinned chapter file opens with \textbackslash{}input\{preamble\} and carries no preamble text of its own, so the project's single Stacks-wide preamble staged at sources/stacks-chow/preamble.tex is excerpted verbatim as three spans, lines 8-12 (the verbatim package together with the project's own reference, slogan and history display-environment definitions, including the author's commented-out older form of the reference environment), lines 14-18 (Xy-pic setup) and lines 154-217 (the native theorem declarations and mathematical macros). The lines 8-12 span is required because a counted statement or its surrounding text may use those display environments, which the project typesets as comments. Nothing was added, removed or redefined; the wrapper class and the amsmath/amssymb/amsthm/graphicx base packages are supplied by the assembler.\par\medskip
\noindent\textit{Editorial note.} Same-chapter pointers are located at their \textbackslash{}label lines inside the pinned chapter file; cross-chapter pointers are located in the project's own tag list staged in the snapshot (sources/stacks-chow/tags.txt), at the verbatim lines carrying their canonical tags.\par\medskip
\noindent\textit{Allowed standard prerequisites (editorial).}\begin{itemize}
\item The characterisation of the Chern classes by one relation on the projective bundle. Same chapter chow.tex, label lemma-determine-intersections, invoked at proof lines 7497 and 7524. Gives the characterising relation for the classes obtained by capping with the Chern classes, and the uniqueness identifying the classes obtained after applying the Gysin map.
\item The Gysin map factors through rational equivalence, commutes with the cap with a first Chern class and with flat pullback. Same chapter chow.tex, labels lemma-gysin-factors (proof line 7514), lemma-gysin-commutes-cap-c1 (proof line 7516) and lemma-gysin-flat-pullback (proof line 7517). They allow the Gysin map to be applied to the relation and the result rewritten in terms of the Gysin images of the classes.
\end{itemize}\medskip
\ifdefined\hypertarget\hypertarget{sourcepointer-1}{}\fi
\noindent\textit{Source reference.} Original definition 02T8: sources/stacks-chow/pinned-original-source.tex, source lines 5069--5069. Navigation pointer for \textbackslash{}ref\{definition-gysin-homomorphism\} used inside the counted span (line 7490); the argument that consumes it is retained in the span.\par\medskip
\ifdefined\hypertarget\hypertarget{sourcepointer-2}{}\fi
\noindent\textit{Source reference.} Original lemma 02U6: sources/stacks-chow/pinned-original-source.tex, source lines 7280--7280. Navigation pointer for \textbackslash{}ref\{lemma-determine-intersections\} used inside the counted span (line 7497); the argument that consumes it is retained in the span.\par\medskip
\ifdefined\hypertarget\hypertarget{sourcepointer-3}{}\fi
\noindent\textit{Source reference.} Original lemma 0B72: sources/stacks-chow/pinned-original-source.tex, source lines 5478--5478. Navigation pointer for \textbackslash{}ref\{lemma-gysin-commutes-cap-c1\} used inside the counted span (line 7516); the argument that consumes it is retained in the span.\par\medskip
\ifdefined\hypertarget\hypertarget{sourcepointer-4}{}\fi
\noindent\textit{Source reference.} Original lemma 02TO: sources/stacks-chow/pinned-original-source.tex, source lines 5420--5420. Navigation pointer for \textbackslash{}ref\{lemma-gysin-factors\} used inside the counted span (line 7514); the argument that consumes it is retained in the span.\par\medskip
\ifdefined\hypertarget\hypertarget{sourcepointer-5}{}\fi
\noindent\textit{Source reference.} Original lemma 0B71: sources/stacks-chow/pinned-original-source.tex, source lines 5303--5303. Navigation pointer for \textbackslash{}ref\{lemma-gysin-flat-pullback\} used inside the counted span (line 7517); the argument that consumes it is retained in the span.\par\medskip
\ifdefined\hypertarget\hypertarget{sourcepointer-6}{}\fi
\noindent\textit{Source reference.} Original reference 02QL: sources/stacks-chow/pinned-original-source.tex, source lines 1372--1372. Navigation pointer for \textbackslash{}ref\{situation-setup\} used inside the counted span (line 7486); the argument that consumes it is retained in the span.\par\medskip
\begin{lemma}
\label{lemma-cap-chern-class-commutes-with-gysin}
Let $(S, \delta)$ be as in Situation \ref{situation-setup}.
Let $X$ be locally of finite type over $S$.
Let $\mathcal{E}$ be a finite locally free sheaf of rank $r$ on $X$.
Let $(\mathcal{L}, s, i : D \to X)$ be as in
Definition \ref{definition-gysin-homomorphism}.
Then $c_j(\mathcal{E}|_D) \cap i^*\alpha = i^*(c_j(\mathcal{E}) \cap \alpha)$
for all $\alpha \in \CH_k(X)$.
\end{lemma}

\begin{proof}
Write $\alpha_j = c_j(\mathcal{E}) \cap \alpha$, so $\alpha_0 = \alpha$.
By Lemma \ref{lemma-determine-intersections} we have
$$
\sum\nolimits_{i = 0}^r
(-1)^i c_1(\mathcal{O}_P(1))^i \cap
\pi^*(\alpha_{r - i}) = 0
$$
in the chow group of the projective bundle
$(\pi : P \to X, \mathcal{O}_P(1))$
associated to $\mathcal{E}$. Consider the fibre product diagram
$$
\xymatrix{
P_D = \mathbf{P}(\mathcal{E}|_D) \ar[r]_-{i'} \ar[d]_{\pi_D} &
P \ar[d]^\pi \\
D \ar[r]^i & X
}
$$
Note that $\mathcal{O}_{P_D}(1)$ is the pullback of $\mathcal{O}_P(1)$.
Apply the Gysin map $(i')^*$ (Lemma \ref{lemma-gysin-factors}) to the
displayed equation above.
Applying Lemmas \ref{lemma-gysin-commutes-cap-c1} and
\ref{lemma-gysin-flat-pullback} we obtain
$$
\sum\nolimits_{i = 0}^r
(-1)^i c_1(\mathcal{O}_{P_D}(1))^i \cap
\pi_D^*(i^*\alpha_{r - i}) = 0
$$
in the chow group of $P_D$.
By the characterization of Lemma \ref{lemma-determine-intersections}
we conclude.
\end{proof}

\end{document}
